\documentclass[professionalfont]{beamer}
%\usecolortheme{default}
\usetheme{Madrid}
\usepackage{subcaption}
\usepackage{caption}
\usepackage{tikz}
\usepackage{dsfont}

\usepackage{pst-all}
\usepackage{leftidx}

\title{\href{https://doi.org/10.21468/SciPostPhys.16.4.100}{\textbf{Symmetries and Anomalies of Kitaev Spin-$S$ Models: Identifying Symmetry-Enforced Exotic Quantum Matter}}}
\subtitle{\href{https://doi.org/10.21468/SciPostPhys.16.4.100}{R. Liu, H. T. Lam, H. Ma, \& L. Zou \\ SciPost Phys. 16, 100 (2024)}}

\author{Chen-Chih Wang}
\institute{NTHU}
\date{\today}

\newcommand{\anglelr}[1]{\left\langle #1 \right\rangle}
\newcommand{\abs}[1]{\left| #1 \right|}
\newcommand{\de}{\partial}
\newcommand{\pdif}[1]{\frac{\partial }{\partial #1}}
\newcommand{\pdiff}[2]{\frac{\partial #1}{\partial #2}}
\newcommand{\lr}[1]{\left(#1\right)}
\newcommand{\lrm}[1]{\left[#1\right]}
\newcommand{\lrg}[1]{\left\{#1\right\}}
\newcommand{\lran}[1]{\left\langle #1 \right\rangle}
\newcommand{\lra}[1]{\left|#1\right|}
\newcommand{\ket}[1]{\left| #1 \right\rangle}
\newcommand{\bra}[1]{\left\langle #1 \right|}
\newcommand{\braket}[2]{\left\langle #1 \right|\left. #2\right\rangle}
\newcommand{\mrm}[1]{\mathrm{#1}}
\newcommand{\bo}[1]{\mathbf{#1}}
\newcommand{\gk}{\gamma_{\mathbf{k}}}
\newcommand{\uk}{u_{\mathbf{k}}}
\newcommand{\vk}{v_{\mathbf{k}}}
\newcommand{\vkp}{v_{\mathbf{k}'}}
\newcommand{\hereref}[2]{\href{#1}{{\tiny (#2)}}}
\newcommand{\tinyeq}[1]{{\tiny\[#1\]}}

\newtheorem{axiom}{Axiom}

\newcommand{\wordfig}[2]{\begin{minipage}[h]{#1 cm}
	\vspace{0pt}
	\includegraphics[width=\textwidth]{#2}
   \end{minipage}}
   
\newcommand{\wordfigdir}[2]{
\begin{minipage}[h]{#1 cm}
\centering
	\vspace{0pt}
	#2
\end{minipage}}

%\setbeamercolor{frametitle}{fg=black,bg=blue}
%\setbeamertemplate{frametitle}[default][left]
%\setbeamerfont{frametitle}{size=\huge}
%\setbeamerfont{block title}{size=\Large}

\begin{document}
\footnotesize
%%%%%%%%%%%%%%%%%%%%%%%%%%%%%%%%%%%%%%%%%%%%%%%%
\begin{frame}
    \titlepage
\end{frame}
%%%%%%%%%%%%%%%%%%%%%%%%%%%%%%%%%%%%%%%%%%%%%%%%
\section{Motivations}
\begin{frame}{Motivations}
    \textbf{\underline{Model:}} 
    \[
        H = -\sum_{\anglelr{ij}}J_{\alpha_{ij}}S^{\alpha_{ij}}_iS^{\alpha_{ij}}_j
    \]
    Useful property
    \[
    e^{i\pi S^\mu}e^{i\pi S^\nu} = \epsilon_{\mu\nu\lambda}e^{i\pi S^\lambda} \quad,\quad e^{i\pi S^\mu}e^{i\pi S^\nu} = (-1)^{2S}e^{i\pi S^\nu}e^{i\pi S^\mu}
    \]
    Conserved $\mathbb{Z}_2$ flux:
    \[
        W_p = \prod_{j\in\de p}e^{i\pi S^{\hat{n}_j}_j} \quad \wordfig{1.5}{honeycomb.pdf}
    \]
    \textbf{\underline{Observations:}} Even-odd effect \hereref{https://doi.org/10.1103/PhysRevLett.130.156701}{H. Ma, PRL 130, 156701 (2023)}
    \begin{itemize}
        \item $S\in\mathbb{Z}+\frac{1}{2}$: gauge charge is \textbf{fermionic} and the $\mathbb{Z}_2$ gauge field is \textbf{deconfined}.
        \item $S\in\mathbb{Z}$: gauge charge is \textbf{bosonic} and the $\mathbb{Z}_2$ gauge field can be Higgsed.
    \end{itemize}
    
    \textbf{\underline{Questions:}}
    \begin{itemize}
        \item To what extent is the even-odd effect stable against perturbation?
        \item Can the results be obtained via a more elementary, direct approach (via spins instead of partons)?
    \end{itemize}
\end{frame}
%%%%%%%%%%%%%%%%%%%%%%%%%%%%%%%%%%%%%%%%%%%%%%%%
\begin{frame}{Symmetries}
    \begin{itemize}
        \item \underline{(Global) $\mathbb{Z}_2^{[1]}$ 1-form symmetry:} Generators are two global loops (on a torus)
        \[
        \prod_{j\in\mathcal{L}_v}e^{i\pi S^{\hat{n}_j}_j} \quad,\quad \prod_{j\in\mathcal{L}_h}e^{i\pi S^{\hat{n}_j}_j}
        \]

        \item \underline{$\mathbb{Z}^T_2$ time reversal symmetry:} $S^\mu_i \to -S^\mu_i$

        \item \underline{$\mathbb{Z}^{(x)}_2 \times \mathbb{Z}^{(y)}_2$}: Generators are
        \[
        \prod_j e^{i\pi S^x_j} \quad,\quad \prod_j e^{i\pi S^y_j}
        \]
        Notice that $e^{i\pi S^x_j}e^{i\pi S^y_j} = e^{i\pi S^z_j}$
    \end{itemize}
    \textbf{Full internal symmetry}: $\mathbb{Z}_2^{[1]} \times \mathbb{Z}^T_2 \times \mathbb{Z}^{(x)}_2 \times \mathbb{Z}^{(y)}_2$
\end{frame}
%%%%%%%%%%%%%%%%%%%%%%%%%%%%%%%%%%%%%%%%%%%%%%%%
\section{Anomalous 1-Form Symmetry}
\begin{frame}{Anomalous 1-Form Symmetry}
    \textbf{\underline{Anomaly}: Presence of obstruction to gauging symmetries.} Here we discuss the 1-form symmetry $\mathbb{Z}_2^{[1]}$.
    \begin{itemize}
        \item The action of a generator of $\mathbb{Z}^{[1]}_2$ symmetry 
        \tinyeq{
            S^\mu_k \to \lr{\prod_{j\in\mathcal{L}}e^{i\pi S^{\hat{n}_j}_j}} S^\mu_k \lr{\prod_{j\in\mathcal{L}}e^{i\pi S^{\hat{n}_j}_j}}^\dagger = \lambda^\mu_kS^\mu_k
        }
        where
        \tinyeq{
            \lambda^\mu_k = 
            \begin{cases}
                -1, & k\in\mathcal{L} ,\;\hat{n}_k\in\mathcal{L} \\
                1 , & \text{otherwise.}
            \end{cases}
        }
        \item Introduce the $\mathbb{Z}_2$ gauge field $A_{ij}$ obey the rule of  gauge transformation: $A_{ij} \to \lambda_i^{\alpha_{ij}} A_{ij} \lambda^{\alpha_{ij}}_j$, and then insert them in between matter fields
        \tinyeq{
             H = -\sum_{\anglelr{ij}}J_{\alpha_{ij}}S^{\alpha_{ij}}_iS^{\alpha_{ij}}_j \to -\sum_{\anglelr{ij}}J_{\alpha_{ij}}S^{\alpha_{ij}}_i A_{ij} S^{\alpha_{ij}}_j
        }
        \item Define the electric fields $E_i^\mu$
        \tinyeq{
            \begin{array}{ccccccc}
                \lr{E_i^\mu}^2=\mathds{1} &,& E_i^{\alpha_{ij}} A_{ij} E_i^{\alpha_{ij}} = - A_{ij} &,& [E_i^\mu, E^\nu_j] = 0 &,& [E_i^\mu, S^\nu_j] = 0
            \end{array}
        }
        
    \end{itemize}
\end{frame}
%%%%%%%%%%%%%%%%%%%%%%%%%%%%%%%%%%%%%%%%%%%%%%%%
\begin{frame}{Anomalous 1-Form Symmetry}
    \begin{itemize}
        \item \item Define the Gauss law (generators of gauge symmetries)
    \end{itemize}
\end{frame}
%%%%%%%%%%%%%%%%%%%%%%%%%%%%%%%%%%%%%%%%%%%%%%%%

%%%%%%%%%%%%%%%%%%%%%%%%%%%%%%%%%%%%%%%%%%%%%%%%

%%%%%%%%%%%%%%%%%%%%%%%%%%%%%%%%%%%%%%%%%%%%%%%%

\end{document}